\section{Класс \texttt{std::string\_view}}

\subsection{Ключевое слово \texttt{auto}: назначение, правила применения, ограничения}
\begin{definition}[\texttt{auto}]
\texttt{auto} — это ключевое слово, которое позволяет компилятору вывести тип переменной из инициализатора.
\end{definition}

Основные назначения:
\begin{itemize}
    \item сокращение записи длинных и сложных типов;
    \item уменьшение жёсткой привязки к конкретной реализации типа;
    \item удобство при работе с итераторами, шаблонами и обобщённым кодом.
\end{itemize}

\begin{example}
\begin{verbatim}
auto x = 42;              // int
auto y = 3.14;            // double
auto it = v.begin();      // тип итератора выводится автоматически
\end{verbatim}
\end{example}

Ограничения:
\begin{itemize}
    \item \texttt{auto} должен иметь инициализатор;
    \item он не задаёт тип сам по себе, а только выводит его;
    \item при использовании \texttt{auto} важно понимать, что вывод может убрать ссылки и \texttt{const}, если не указать \texttt{auto\&} или \texttt{const auto\&}.
\end{itemize}

\begin{remark}
На практике \texttt{auto} особенно полезен там, где тип либо слишком длинный, либо несущественен для понимания кода.
\end{remark}

\subsection{Строковые литералы, глобальная доступность и слияние одинаковых литералов}
\begin{definition}[Строковый литерал]
Строковый литерал — это текст в двойных кавычках, например \texttt{"hello"}.
\end{definition}

Строковые литералы:
\begin{itemize}
    \item имеют статическое время жизни;
    \item обычно доступны на всём протяжении выполнения программы;
    \item являются глобально доступными в смысле времени жизни и адресуемости.
\end{itemize}

\begin{remark}
Одинаковые строковые литералы компилятор или линкер может \textit{сливать} (\textit{merge}): несколько одинаковых литералов могут указывать на один и тот же участок памяти. Но полагаться на это как на строгое правило языка не следует.
\end{remark}

\subsection{\texttt{std::string\_view} как разновидность косвенной адресации}
\begin{definition}[std::string\_view]
\texttt{std::string\_view} — это невладеющее представление последовательности символов: он хранит ссылку на данные и длину, но не управляет памятью.
\end{definition}

С точки зрения косвенной адресации это:
\begin{itemize}
    \item не копия строки;
    \item не владелец данных;
    \item лёгкая обёртка над «окном» в строку;
    \item средство доступа к подстроке без выделения новой памяти.
\end{itemize}

\subsection{Какие объекты нужны для представления Null-terminated строки?}
Для представления Null-terminated строки обычно достаточно:
\begin{itemize}
    \item одного массива символов \texttt{char[]};
    \item или одного указателя \texttt{char*}, если сам массив уже существует и конец строки обозначен символом \texttt{'\textbackslash 0'}.
\end{itemize}

\begin{example}
\begin{verbatim}
char s[] = "hello";
const char* p = s;
\end{verbatim}
\end{example}

\begin{remark}
Смысл Null-terminated строки задаётся именно наличием завершающего нулевого символа, а не отдельным объектом длины.
\end{remark}

\subsection{Какие объекты нужны для представления суффикса, префикса и инфикса Null-terminated строки?}
Для подстрок возможны разные представления.

\begin{itemize}
    \item \textbf{Префикс}: достаточно указателя на начало строки, если подстрока заканчивается \texttt{\textbackslash0}, либо указателя + длина, если конец не совпадает с физическим концом.
    \item \textbf{Суффикс}: часто достаточно указателя на начало суффикса, если строка уже null-terminated.
    \item \textbf{Инфикс}: обычно нужны либо два указателя \texttt{[first, last)}, либо указатель + длина.
\end{itemize}

\begin{remark}
Для общего случая подстроки удобнее хранить не только начало, но и длину. Тогда можно обращаться и к суффиксу, и к инфиксу без поиска конца.
\end{remark}

\subsection{Два указателя против указателя + длина}
Есть два основных способа представить подстроку:
\begin{itemize}
    \item \textbf{пара указателей} \texttt{[first, last)};
    \item \textbf{указатель + длина} \texttt{(ptr, size)}.
\end{itemize}

\textbf{Пара указателей} хорошо соотносится с концепцией range в C++:
\[
[first, last)
\]
где \texttt{last} указывает на элемент \textbf{после} последнего.

Плюсы:
\begin{itemize}
    \item естественная модель диапазона;
    \item удобна для алгоритмов;
    \item не нужно отдельно хранить длину.
\end{itemize}

Минусы:
\begin{itemize}
    \item длину приходится вычислять как расстояние между указателями;
    \item не всякий код удобно писать в таком виде.
\end{itemize}

\textbf{Указатель + длина}:
\begin{itemize}
    \item удобно хранить длину явно;
    \item быстро проверять размер;
    \item удобно для строковых view.
\end{itemize}

Минусы:
\begin{itemize}
    \item нужно следить за согласованностью начала и длины;
    \item это не совсем та же форма, что классический range \texttt{[first, last)}.
\end{itemize}

\subsection{Передача и возврат \texttt{std::string\_view} из функций}
\texttt{std::string\_view} обычно передают и возвращают \textbf{по значению}.

\begin{example}
\begin{verbatim}
void print(std::string_view sv) {
    std::cout << sv << '\n';
}

std::string_view head(std::string_view sv) {
    return sv.substr(0, 3);
}
\end{verbatim}
\end{example}

\begin{remark}
Такой объект считается «нетяжёлым»: он обычно хранит только указатель и длину, поэтому копировать его дёшево.
\end{remark}

\subsection{Конструирование и копирование \texttt{std::string\_view}. Почему это невладеющий тип?}
\texttt{std::string\_view} можно создать из:
\begin{itemize}
    \item строкового литерала;
    \item C-строки;
    \item \texttt{std::string};
    \item другого \texttt{std::string\_view}.
\end{itemize}

\begin{example}
\begin{verbatim}
std::string s = "hello";
std::string_view v1 = s;
std::string_view v2 = "world";
std::string_view v3 = v1;
\end{verbatim}
\end{example}

Он \textbf{не владеет} данными:
\begin{itemize}
    \item не выделяет память под строку;
    \item не освобождает память;
    \item лишь ссылается на уже существующий буфер.
\end{itemize}

\begin{remark}
Из-за этого главное ограничение \texttt{string\_view}: он не должен жить дольше, чем данные, на которые ссылается.
\end{remark}

\subsection{Ограничения \texttt{std::string\_view} как невладеющего типа}
Поскольку \texttt{string\_view} не владеет данными, возникают ограничения:
\begin{itemize}
    \item нельзя сохранять view дольше, чем живёт исходная строка;
    \item нельзя использовать его как замену владению;
    \item нужно быть осторожным с временными объектами;
    \item модификация исходной строки может сделать view невалидным.
\end{itemize}

\begin{example}
\begin{verbatim}
std::string_view bad() {
    std::string s = "temp";
    return s; // dangling view
}
\end{verbatim}
\end{example}

\subsection{Ограничение \texttt{std::string\_view} при работе со стандартной библиотекой}
\texttt{std::string\_view} не всегда можно передавать туда, где ждут \texttt{const char*}.

\begin{itemize}
    \item если функция принимает именно C-строку, ей нужен нуль-терминированный буфер;
    \item \texttt{string\_view} гарантирует диапазон символов, но не всегда гарантирует наличие завершающего \texttt{'\textbackslash0'} в конце именно подстроки;
    \item поэтому иногда требуется явное преобразование в \texttt{std::string} или использование \texttt{.data()} с осторожностью.
\end{itemize}

\begin{remark}
Если API библиотеки требует \texttt{const char*} и предполагает нуль-терминацию, \texttt{string\_view} можно использовать не во всех случаях без дополнительной проверки.
\end{remark}

\subsection{Ключевые операции над строками с помощью \texttt{std::string\_view}}
Основные операции:
\begin{itemize}
    \item \texttt{size()}, \texttt{length()};
    \item \texttt{empty()};
    \item \texttt{substr()};
    \item \texttt{remove\_prefix()};
    \item \texttt{remove\_suffix()};
    \item \texttt{compare()};
    \item \texttt{starts\_with()}, \texttt{ends\_with()} в современных стандартах.
\end{itemize}

\begin{example}
\begin{verbatim}
std::string_view sv = "hello world";
auto prefix = sv.substr(0, 5);
sv.remove_prefix(6);
\end{verbatim}
\end{example}

\subsection{Что происходит, если модифицировать подстроку через \texttt{std::string\_view}?}
\texttt{std::string\_view} сам по себе не модифицирует данные, потому что он не владеет памятью и не является изменяемой строкой.

\begin{itemize}
    \item можно изменить сам \textbf{view}: например, сдвинуть начало через \texttt{remove\_prefix()};
    \item нельзя изменить символы исходной строки через \texttt{string\_view};
    \item если нужна модификация текста, нужно работать со \texttt{std::string} или другим изменяемым буфером.
\end{itemize}

\begin{remark}
После изменения самого view это уже тот же объект \texttt{string\_view}, но он начинает указывать на другой участок того же исходного буфера.
\end{remark}

\subsection{Краткие ответы на контрольные вопросы}
\begin{itemize}
    \item \textbf{Какие объекты нужны для null-terminated строки?} Обычно один массив \texttt{char[]} или указатель на него.
    \item \textbf{Какие объекты нужны для подстроки?} Для общего случая удобны либо пара указателей, либо указатель + длина.
    \item \textbf{Как связаны пары \texttt{char*} и \texttt{char* + size\_t}?} Это два способа задать границы подстроки; они не полностью взаимозаменяемы без потери информации о конце или длине.
    \item \textbf{Как передают \texttt{string\_view}?} Обычно по значению.
    \item \textbf{Глобально ли доступен строковый литерал?} Да, по времени жизни он статический.
    \item \textbf{Гарантирует ли \texttt{string\_view} нуль-терминацию?} Нет, только если view охватывает строку целиком и источник это обеспечивает.
    \item \textbf{Можно ли всегда подставить \texttt{string\_view} вместо \texttt{const char*}?} Нет, только когда требуется нестрогая работа с диапазоном или когда гарантирована нуль-терминация.
    \item \textbf{Можно ли модифицировать подстроку через \texttt{string\_view}?} Сам view — да (например, сдвиг границ), сами символы — нет.
\end{itemize}
